\section{Received Traditions Examined}

A saint who governed the Western imperial capital for twenty-three years inevitably gathered a dense halo of stories, liturgical attributions, and dramatic set pieces. Respect for Christian memory requires neither credulity nor mockery: each memorable tradition is tested against the earliest attested sources so that historical fact, fifth-century hagiographic shaping, medieval legend, and liturgical growth stand in their true colors.

\subsection{The Swarm of Bees on the Infant's Mouth}

Iconography routinely depicts Saint Ambrose with a beehive, honoring him as the ``honey-tongued Doctor'' whose eloquence nourished the Church. The motif traces to the third chapter of Paulinus of Milan's \work{Vita Ambrosii} (c.~422), which relates that while the infant Ambrose slept with open mouth in the courtyard of the praetorian residence in Gaul, a swarm of bees settled over his face, entered and left his mouth without stinging him, and rose out of sight into the sky---a sign his father took to foretell future greatness.

No contemporary writing of Ambrose, Augustine, or Basil mentions this childhood prodigy. Moreover, the swarm of bees settling on an infant's lips is a well-known classical literary \latin{topos} applied in antiquity to eloquent poets and philosophers such as Pindar and Plato. As biography, the scene rests on a single fifth-century hagiographic witness employing an inherited classical omen, while in iconography it aptly symbolizes the enduring sweetness of Ambrose's preaching and hymns.

\subsection{The Child's Cry: ``Ambrose, Bishop!''}

According to Paulinus of Milan (summarized in de Romestin's Prolegomena~\S V and notes on \work{De Officiis} I.1.2), when Catholics and Arians clashed in the Milanese basilica after Auxentius died in 374, a child's voice suddenly rang out, ``Ambrose, bishop!'', immediately uniting the divided crowd in acclamation; Ambrose then tried to repel the voters by ordering torture on the bench and fleeing Milan by night until Valentinian~I confirmed his election.

Contemporary sources establish the decisive historical core without those narrative details. Ambrose himself confirms that the Milanese people demanded him from Valentinian~I (\work{Letter} XXI.7), that he tried to avoid the office (\work{De Officiis} I.1.2), and that he was carried straight from the tribunal to the priesthood while still unbaptized (\work{De Officiis} I.1.4). St.~Basil's contemporary \work{Letter} CXCVII.1 corroborates the sudden transfer of a prominent imperial governor to the Milanese see. The eleventh-century \work{Manuale Ambrosianum} later wove the acclamation into Milanese liturgical song---echoing \latin{Ambrosium episcopum} in the Vespers hymn \latin{Miraculum laudabile} (Pars Altera, p.~18) and the child's voice (\latin{Vox infantis sonuit: Dignus Ambrosius, hic erit episcopus}) in morning antiphon \latin{Ant.\ xxiii} (pp.~21--22). The child's cry and Ambrose's theatrical disqualifications belong specifically to Paulinus's fifth-century hagiographic narrative (summarized by de Romestin) and cannot be independently verified from fourth-century testimony.

\subsection{Alternating Composition of the \work{Te Deum} at Easter 387}

A famous medieval story relates that as Ambrose drew Augustine from the Milanese baptismal font at Easter 387, the two Doctors burst into inspired song, improvising alternate verses of the \work{Te Deum laudamus} (``We praise thee, O God''). J.~G. Pilkington's note on \work{Confessions} IX.6.14 in NPNF~I.1 records that English liturgical offices long called the canticle ``The Song of Ambrose and Augustin.''

Contemporary testimony contradicts the story by its complete silence. Augustine devotes \work{Confessions} IX.6--7 to the exact emotional memory of his Easter 387 baptism and the weeping joy of Milan's hymns---and later quotes Ambrose's evening hymn \work{Deus creator omnium} verbatim in IX.12.32---without a word about improvising the \work{Te Deum}. Neither Ambrose's letters nor Paulinus's biography mentions any such miracle. In modern hymnological consensus, the alternating-composition legend first surfaces centuries later in the early Middle Ages, while critical scholarship commonly associates the Latin text of the \work{Te Deum} with Bishop Nicetas of Remesiana (late fourth or early fifth century). The canticle remains a treasured ``Ambrosian hymn'' in Catholic liturgical tradition, not a historical transcript of Easter 387.

\subsection{Barring Emperor Theodosius at the Cathedral Doors}

Few scenes in Christian art are more familiar than Ambrose standing vested at the bronze doors of the Milanese cathedral, hand raised to halt Emperor Theodosius~I and his armed retinue after the massacre at Thessalonica in 390.

Ambrose's own contemporary \work{Letter} LI tells a subtler, more pastoral story. Far from waiting on the cathedral steps for a public showdown, Ambrose left Milan before Theodosius arrived---citing severe bodily sickness and his deliberate wish not to inflame the emperor's volatile temper in public (\work{Letter} LI.5). Writing entirely in his own hand ``that which you alone may read,'' he notified Theodosius privately that he dared not celebrate the Eucharist in the emperor's presence until Theodosius followed King David in tears and penitence (\work{Letter} LI.13--14). What did happen publicly was Ambrose's earlier altar-side confrontation with Theodosius over Callinicum in 388 (\work{Letter} XLI.27--28) and Theodosius's subsequent public penance after Thessalonica. The vivid doorway confrontation enters history through the fifth-century Eastern historian Theodoret of Cyrrhus (\work{Historia Ecclesiastica} V.18), who dramatized the authentic moral triumph of episcopal discipline into an unforgettable tableau.

\subsection{Authorship of the Ambrosian Rite and the ``Ambrosian'' Hymnal}

Popular shorthand often speaks as though Ambrose single-handedly composed the entire Ambrosian Rite of Milan and every hymn called ``Ambrosian.'' Both claims overstate the evidence, as Milan's own liturgical books candidly recognize.

As Cardinal Cesare Monti observed in the preface to the 1640 \work{Missale Ambrosianum}, Ambrose ``not only retained the rites already instituted'' by Milanese predecessors such as Myrocles, ``but increased them with such profit and fruitfulness that not long after his time our Church received the distinguished name of Ambrosian,'' after which Simplicianus, Eusebius, and later archbishops continued to add and restore liturgical texts. Ambrose's verifiable liturgical legacy is substantial and specific: the introduction of Eastern-style antiphonal psalmody and congregational hymnody (\work{Confessions} IX.7.15; \work{Sermon against Auxentius} 34), the post-baptismal catechesis of \work{De Mysteriis}, and the Milanese Eucharistic prayers witnessed alongside \work{De Sacramentis}. The fully developed medieval Ambrosian Missal and \work{Manuale} represent centuries of organic local evolution around that fourth-century core.

Similarly, while scores of Latin hymns bear the title \latin{Ambrosianus} in medieval manuscripts and monastic rules (where the adjective often simply denotes the iambic dimeter genre), external fourth- and fifth-century testimony securely authenticates four core hymns---namely \work{Aeterne rerum conditor}, \work{Deus creator omnium}, \work{Iam surgit hora tertia}, and \work{Intende qui regis Israel} / \work{Veni Redemptor gentium}---with a small circle of further candidates (such as \work{Illuminans Altissimus}, cited by Cassiodorus) accepted by Maurist editors (\work{Prolegomena}~\S VI).

\subsection{Tomb, Relics, and the 1864 Basilica Excavations}

Ambrose's own \work{Letter} XXII.13 proves that in June 386 he relinquished the burial space beneath the right side of the altar in his new Basilica Ambrosiana to Saints Gervasius and Protasius, intending to rest beside them; he was laid there on Easter Sunday, 5~April 397. According to de Romestin's Prolegomena~\S II, in 835 Archbishop Angilbert~II of Milan placed Ambrose's remains together with the two martyrs in a single great porphyry sarcophagus beneath the high altar, and excavations in 1864 uncovered both that porphyry sarcophagus and the two original marble-lined fourth-century vaults beside it, containing fourth-century coins of Flavius Victor and Theodosius along with fragments of cloth of gold.

Two strict cautions govern that relic record. First, H.~de Romestin's Prolegomena~\S II suffers from a striking editorial slip: it misnames the two martyrs buried beside Ambrose as ``SS.\ Cosmas and Damian, a.d.~389,'' contradicting Ambrose's own \work{Letter} XXII, Augustine's \work{Confessions} IX.7.16, and de Romestin's own chronological table in \S III (which correctly dates the finding of Saints Gervasius and Protasius to 386). Second, the 835 translation and 1864 excavation are known here through nineteenth-century historical summaries rather than primary archaeological excavation reports from the Basilica Ambrosiana. Ambrose's 397 burial beside Gervasius and Protasius in the Basilica Ambrosiana is historically secure; forensic authentication of the shrine's modern relics lies outside historical biography.
